\section*{الفصل \ref{ch:b1:logic} --- المنطق والمجموعات والتطبيقات}

\begin{solution}{exo:b1:logic:1}
النفي، بتمرير $\lnot$ عبر كل مسوِّر
(\cref{prop:b1:logic:negquant}) وباستعمال $\lnot(P \implies Q) \iff
P \land \lnot Q$:
\begin{enumerate}
  \item $\exists x \in \R,\ \forall y \in \R,\ x + y \leq 0$؛
  \item $\forall x \in \R,\ \exists y \in \R,\ xy \neq 0$؛
  \item $\exists \varepsilon > 0,\ \forall \delta > 0,\ \exists x \in
        \R,\ \abs{x} \leq \delta \text{ و} \abs{f(x)} >
        \varepsilon$.
\end{enumerate}
\omterm{def:b1:logic:statement}{العبارة} (1) صادقة: من أجل $x$ معطى، خذ $y = -x + 1$؛ عندئذ $x + y = 1 >
0$. \omterm{def:b1:logic:statement}{والعبارة} (2) صادقة: العدد $x = 0$ يحقق $xy = 0$ لكل $y$.
\end{solution}

\begin{solution}{exo:b1:logic:2}
جدول الحقيقة، بكتابة ص/ك في الحالات الأربع $(P, Q)$:
\begin{center}
\begin{tabular}{cc|c|c|c|c}
$P$ & $Q$ & $P \implies Q$ & $\lnot(P \implies Q)$ & $\lnot Q$ &
$P \land \lnot Q$ \\
\hline
ص & ص & ص & ك & ك & ك \\
ص & ك & ك & ص & ص & ص \\
ك & ص & ص & ك & ك & ك \\
ك & ك & ص & ك & ص & ك \\
\end{tabular}
\end{center}
يتطابق العمودان $4$ و $6$، وهذا يبرهن على التكافؤ. ومنه فإن نفي
\omterm{def:b1:logic:statement}{عبارة} «إذا كانت دالة قابلة للاشتقاق فهي متصلة» هو:
«توجد دالة قابلة للاشتقاق وغير متصلة»
(وهي \omterm{def:b1:logic:statement}{عبارة} كاذبة في واقع الأمر، إذ إن الاستلزام الأصلي صادق،
انظر \cref{ch:b1:derivative}).
\end{solution}

\begin{solution}{exo:b1:logic:3}
\emph{\omterm{prop:b1:logic:rules}{بعكس النقيض}.} نفترض أن $x < 1$. عندئذ $x^3 < 1$ (دالة التكعيب
متزايدة) و $x < 1$، ومنه $x^3 + x < 2$. وهذا يبرهن على \omterm{prop:b1:logic:rules}{عكس النقيض}،
ومنه على \omterm{def:b1:logic:statement}{العبارة} نفسها.

\emph{بالخلف.} نفترض أن $a > 0$ أصغر عدد حقيقي موجب تمامًا. عندئذ
$a/2$ موجب تمامًا و $a/2 < a$ (لأن $a > 0$)، وهذا يناقض الأصغرية.
ومنه لا وجود لعدد $a$ كهذا.
\end{solution}

\begin{solution}{exo:b1:logic:4}
\begin{enumerate}
  \item حالة البداية $n = 0$: $2^0 = 1 = 2^1 - 1$. خطوة الانتقال:
        بافتراض المتطابقة من أجل $n$،
        \[
        \sum_{k=0}^{n+1} 2^k = (2^{n+1} - 1) + 2^{n+1}
        = 2 \cdot 2^{n+1} - 1 = 2^{n+2} - 1 .
        \]
  \item حالة البداية $n = 0$: $4^0 + 5 = 6 = 3 \times 2$. خطوة الانتقال: إذا كان
        $4^n + 5 = 3m$ فإن
        \[
        4^{n+1} + 5 = 4(4^n + 5) - 15 = 3(4m - 5),
        \]
        وهو يقبل القسمة على $3$.
\end{enumerate}
\end{solution}

\begin{solution}{exo:b1:logic:5}
تفترض خطوة الانتقال ضمنًا أن المجموعتين («الأقلام $n$ الأولى»
و«الأقلام $n$ الأخيرة») تتقاطعان، فتنقل الأقلام المشتركة اللون من
إحداهما إلى الأخرى. ومن أجل $n + 1 = 2$ تكون المجموعتان
$\{$القلم الأول$\}$ و $\{$القلم الثاني$\}$: وهما منفصلتان،
فتنهار الحجة. إذن لم يُبرهن قط على $P(1) \implies P(2)$، وينهار
الاستقراء --- وإن كان $P(n) \implies P(n+1)$ صالحًا لكل
$n \geq 2$.
\end{solution}

\begin{solution}{exo:b1:logic:6}
\begin{enumerate}
  \item $x \in A \setminus B \iff x \in A \land x \notin B \iff x \in
        A \land x \in \overline{B} \iff x \in A \cap \overline{B}$.
  \item باستعمال (1) والتوزيعية
        (\cref{prop:b1:logic:setalgebra}):
        $(A \cup B) \cap \overline{C} = (A \cap \overline{C}) \cup
        (B \cap \overline{C})$.
  \item نفترض $A \subseteq B$. عندئذ $A \cup B \subseteq B$ (فكلتا
        القطعتين تقع في $B$) و $B \subseteq A \cup B$ دائمًا، ومنه
        $A \cup B = B$. ونفترض $A \cup B = B$: عندئذ $A \cap B
        \subseteq A$ دائمًا، ويعطي $A \subseteq A \cup B = B$ أن
        $A \subseteq A \cap B$، ومنه $A \cap B = A$. ونفترض $A \cap B =
        A$: عندئذ $A = A \cap B \subseteq B$. فالشروط الثلاثة
        متكافئة إذن (فقد برهنّا على دورة من الاستلزامات).
\end{enumerate}
\end{solution}

\begin{solution}{exo:b1:logic:7}
\begin{enumerate}
  \item \omterm{def:b1:logic:map}{التطبيق} $f$ متباين ($n + 1 = m + 1 \implies n = m$) لكنه غير
        شامل: إذ ليس للعدد $0$ سابقة في $\N$.
  \item \omterm{def:b1:logic:map}{التطبيق} $g$ تقابليّ: \omterm{def:b1:logic:map}{فالتطبيق} $n \mapsto n - 1$ مقلوب ثنائي
        الجانب له على $\Z$.
  \item \omterm{def:b1:logic:map}{التطبيق} $h$ متباين: إذ يعطي $\frac{x+1}{x-1} = \frac{x'+1}{x'-1}$ أن
        $(x+1)(x'-1) = (x'+1)(x-1)$، أي $xx' - x + x' - 1 = xx' -
        x' + x - 1$، ومنه $2x' = 2x$. وهو غير شامل على $\R$:
        فحلّ المعادلة $y = \frac{x+1}{x-1}$ يعطي $x(y - 1) = y + 1$،
        وهي بلا حلّ عندما $y = 1$ (إذ تصير المعادلة $0 = 2$).
        وبأخذ \omterm{def:b1:logic:sets}{مجموعة} الوصول $\R \setminus \{1\}$ يعطي الحساب نفسه
        السابقة الوحيدة $x = \frac{y+1}{y-1}$، ومنه فإن
        $h \colon \R \setminus \{1\} \to \R \setminus \{1\}$
        تقابليّ و $h^{-1}(y) = \frac{y+1}{y-1} = h(y)$: أي أن $h$ مقلوب
        نفسه.
\end{enumerate}
\end{solution}

\begin{solution}{exo:b1:logic:8}
\begin{enumerate}
  \item $x \in f^{-1}(B \cap B') \iff f(x) \in B \cap B' \iff f(x) \in
        B \land f(x) \in B' \iff x \in f^{-1}(B) \cap f^{-1}(B')$.
        وأمّا الصور: فيكون $y \in f(A \cup A')$ إذا وفقط إذا كان
        $y = f(x)$ من أجل $x$ ما في $A$ أو في $A'$، أي إذا وفقط إذا
        كان $y \in f(A)$ أو $y \in f(A')$.
  \item إذا كان $y \in f(A \cap A')$ فإن $y = f(x)$ مع $x \in A$
        و $x \in A'$، ومنه $y \in f(A)$ و $y \in f(A')$. وأمّا
        الاحتواء التام: فخذ $f \colon \R \to \R$، $x \mapsto x^2$،
        $A = \{-1\}$، $A' = \{1\}$: عندئذ $f(A \cap A') = f(\emptyset) = \emptyset$
        بينما $f(A) \cap f(A') = \{1\}$.
  \item ($\Leftarrow$) بأخذ $A = \{x\}$ و $A' = \{x'\}$ حيث $x \neq
        x'$: إذا كان $f(x) = f(x')$ فإن $f(A) \cap f(A') = \{f(x)\}$
        في حين أن $f(A \cap A') = \emptyset$، وهذا يناقض التساوي
        المفترض؛ إذن $f$ متباين. ($\Rightarrow$) ليكن $f$ \omterm{def:b1:logic:inj}{متباينًا}
        وليكن $y \in f(A) \cap f(A')$: أي $y = f(x) = f(x')$ مع
        $x \in A$ و $x' \in A'$؛ ويعطي التباين أن $x = x' \in A \cap
        A'$، ومنه $y \in f(A \cap A')$. ومع (2) يتحقق التساوي.
\end{enumerate}
\end{solution}

\begin{solution}{exo:b1:logic:9}
التركيب $g \circ f = \mathrm{id}_E$ متباين وشامل، فحسب
\cref{prop:b1:logic:comp}~(2) يكون $f$ \omterm{def:b1:logic:inj}{متباينًا} و $g$ \omterm{def:b1:logic:inj}{شاملًا}.
مثال: $E = \N$، $F = \Z$، و $f$ هو الحقن $n \mapsto n$، و
$g \colon \Z \to \N$ حيث $g(n) = n$ من أجل $n \geq 0$ و $g(n) = 0$
من أجل $n < 0$. عندئذ $g(f(n)) = n$ لكل $n \in \N$، لكن $f$ غير
شامل و $g$ غير متباين.
\end{solution}

\begin{solution}{exo:b1:logic:10}
$x^2 - y^2 = x - y \iff (x - y)(x + y) = x - y \iff (x - y)(x + y - 1)
= 0 \iff y = x$ أو $y = 1 - x$. \emph{انعكاسية:} الاختيار $y = x$ يفي بالغرض.
\emph{تماثلية:} الشرط «$y = x$ أو $y = 1 - x$»
متناظر في $x$ و $y$ (إذ إن $y = 1 - x$ يكافئ $x = 1 - y$).
\emph{متعدّية:} نفترض $x \mathbin{\mathcal{R}} y$ و$y
\mathbin{\mathcal{R}} z$؛ وبمراجعة الحالات الأربع نجد أن $z$ يساوي $x$
أو $1 - x$ في كل مرة (مثلًا $y = 1 - x$ و $z = 1 - y$ يعطيان $z = x$).
إذن $\mathcal{R}$ \omterm{def:b1:logic:equiv}{علاقة تكافؤ} و$\mathrm{cl}(x) =
\{x,\, 1 - x\}$. ولهذا الصف عنصر واحد بالضبط عندما يكون $x = 1 - x$،
أي من أجل $x = \frac12$.
\end{solution}

\begin{solution}{exo:b1:logic:11}
ليكن $f \colon E \to \mathcal{P}(E)$ تطبيقًا كيفيًا ولنضع
$D = \{x \in E : x \notin f(x)\} \in \mathcal{P}(E)$. ونفترض $D =
f(a)$ من أجل $a \in E$ ما. إذا كان $a \in D$ فإن $D$ بحكم تعريفه
يعطي $a \notin f(a) = D$: وهذا تناقض. وإذا كان $a \notin D$ فإن $a \notin
f(a)$، ومنه $a \in D$ بحكم تعريف $D$: وهذا تناقض. ومنه فإن $D$
ليست في صورة $f$، فلا يكون $f$ \omterm{def:b1:logic:inj}{شاملًا}. (وبوجه خاص لا توجد \omterm{def:b1:logic:sets}{مجموعة}
في تقابل مع \omterm{def:b1:logic:sets}{مجموعة} أجزائها: فأجزاء $\N$ «أكثر» من الأعداد الصحيحة.)
\end{solution}

\begin{solution}{exo:b1:logic:12}
\begin{enumerate}
  \item ($\Rightarrow$) ليكن $f$ \omterm{def:b1:logic:inj}{شاملًا} وليكن $\Phi(B) =
        \Phi(B')$. من أجل $y \in B$، اختر $x$ يحقق $f(x) = y$؛ عندئذ $x
        \in f^{-1}(B) = f^{-1}(B')$، ومنه $y = f(x) \in B'$. إذن $B
        \subseteq B'$، وبالتناظر $B' \subseteq B$: أي أن $\Phi$
        متباين. ($\Leftarrow$) إذا لم يكن $f$ \omterm{def:b1:logic:inj}{شاملًا} فاختر $y_0
        \in F$ خارج الصورة؛ عندئذ $f^{-1}(\{y_0\}) = \emptyset =
        f^{-1}(\emptyset)$ مع $\{y_0\} \neq \emptyset$، ومنه لا يكون $\Phi$
        \omterm{def:b1:logic:inj}{متباينًا}.
  \item ($\Rightarrow$) ليكن $f$ \omterm{def:b1:logic:inj}{متباينًا} ولتكن $A \subseteq E$.
        نضع $B = f(A)$؛ عندئذ $f^{-1}(B) = \{x : f(x) \in f(A)\}$، ويعطي
        التباين أن $f(x) \in f(A) \iff x \in A$، ومنه $\Phi(B) =
        A$: أي أن $\Phi$ شامل. ($\Leftarrow$) إذا لم يكن $f$
        \omterm{def:b1:logic:inj}{متباينًا} فخذ $x \neq x'$ يحققان $f(x) = f(x')$. عندئذ تحتوي
        كل \omterm{def:b1:logic:sets}{مجموعة} سوابق $f^{-1}(B)$ على $x$ إذا وفقط إذا احتوت
        على $x'$؛ ومنه فإن $\{x\}$ ليست على الصورة $\Phi(B)$، ولا
        يكون $\Phi$ \omterm{def:b1:logic:inj}{شاملًا}.
\end{enumerate}
\end{solution}

\begin{solution}{pb:b1:logic:1}
\textbf{1.} \emph{انعكاسية:} \omterm{def:b1:logic:map}{التطبيق} $\mathrm{id}_E$ تقابل من $E$
على نفسها. \emph{تماثلية:} إذا كان $f \colon E \to F$ \omterm{def:b1:logic:inj}{تقابليًا}
أعطت \cref{thm:b1:logic:inverse} \omterm{def:b1:logic:map}{التطبيق} $f^{-1} \colon F \to E$، وهو نفسه
تقابليّ. \emph{متعدّية:} إذا كان $f \colon E \to F$ و$g \colon F
\to G$ تقابلين قالت \cref{prop:b1:logic:comp}~(1) إن $g \circ f
\colon E \to G$ تقابل. (وهذا «شبيه» \omterm{def:b1:logic:equiv}{بعلاقة تكافؤ} لا غير: إذ إن
\omterm{def:b1:logic:sets}{مجموعة} كل المجموعات ليست هي نفسها \omterm{def:b1:logic:sets}{مجموعة}، بسبب المفارقات التي
يلمّح إليها \cref{exo:b1:logic:11}؛ والمهم هو الخصائص الثلاث.)

\textbf{2.} إذا كان $f \colon E \to F$ و $g \colon F \to G$
متباينين كان $g \circ f$ \omterm{def:b1:logic:inj}{متباينًا} حسب \cref{prop:b1:logic:comp}~(1):
أي $E \preceq G$. وأمّا النقطة الثانية فنضيّق \omterm{def:b1:logic:sets}{مجموعة} وصول $f$ إلى
صورته: \omterm{def:b1:logic:map}{فالتطبيق} $\tilde f \colon E \to f(E)$، $x \mapsto f(x)$، شامل بحكم
إنشاء $f(E)$ ومتباين لأن $f$ متباين، فهو إذن تقابليّ:
أي $E \approx f(E)$.

\textbf{3.} ($\Rightarrow$) ليكن $f \colon E \to F$ \omterm{def:b1:logic:inj}{متباينًا}
ولنثبّت $a \in E$ (إذ $E \neq \emptyset$). نعرّف $s \colon F \to E$
كما يلي: $s(y)$ هو العنصر الوحيد $x$ الذي يحقق $f(x) = y$ عندما
$y \in f(E)$ (والوحدانية بالتباين)، و $s(y) = a$ فيما عدا ذلك.
ومن أجل كل $x
\in E$ لدينا $s(f(x)) = x$، فيُبلغ كل $x$: أي أن $s$ شامل.
($\Leftarrow$) ليكن $s \colon F \to E$ \omterm{def:b1:logic:inj}{شاملًا}. من أجل كل $x \in
E$ اختر $y_x \in F$ واحدًا يحقق $s(y_x) = x$، وضع $u(x) = y_x$. إذا
كان $u(x) = u(x')$ فإن $x = s(u(x)) = s(u(x')) = x'$: أي أن $u \colon E \to F$
متباين.

\textbf{4.} \omterm{def:b1:logic:map}{التطبيق} $n \mapsto n + 1$ يرسل $\N$ في $\N^*$، وهو
متباين ($n + 1 = m + 1 \implies n = m$) وشامل (إذ إن كل $m \geq 1$ يساوي
$(m - 1) + 1$ مع $m - 1 \in \N$). وأمّا $\sigma$: فهو يرسل الأعداد
الزوجية $0, 2, 4, \dots$ إلى $0, 1, 2, \dots$ والأعداد الفردية $1, 3,
5, \dots$ إلى $-1, -2, -3, \dots$ والتباين: الدخائل الزوجية تقع في
$\N$ ($\sigma(n) = n/2 \geq 0$) والدخائل الفردية تقع في الأعداد
الصحيحة السالبة تمامًا ($\sigma(n) = -(n+1)/2 \leq -1$)، فأيّ
تصادم لا بد أن يقع داخل صف زوجية واحد، حيث يكون $\sigma$ رتيبًا
تمامًا ($n/2 = m/2$ أو $(n+1)/2 = (m+1)/2$ يفرضان $n =
m$). والشمول: كل $k \geq 0$ يساوي $\sigma(2k)$؛ وكل $k \leq -1$
يساوي $\sigma(-2k - 1)$ مع $-2k - 1 \geq 1$ فرديًا. إذن
$\N \approx \N^*$ و $\N \approx \Z$.

\textbf{5.} لدينا $C_0 = E \setminus g(F) \subseteq C$، ومنه فإن $x \notin C$
يستلزم $x \notin C_0$، أي $x \in g(F)$: أي أن عنصرًا $y \in F$ يحقق
$g(y) = x$. وإذا كان $g(y') = x$ أيضًا أعطى تباين $g$ أن $y' = y$.
ومنه فإن الشرط الثاني من تعريف $h$ يعيّن عنصرًا وحيدًا معرَّفًا
تعريفًا سليمًا هو $g^{-1}(x)$.

\textbf{6.} الصور المباشرة تتبادل مع الاتحادات
(\cref{exo:b1:logic:8}~(1)، مطبَّقة على $f$ ثم على $g$):
\[
g\bigl(f(C)\bigr)
= g\Bigl(f\Bigl(\bigcup_{n \in \N} C_n\Bigr)\Bigr)
= \bigcup_{n \in \N} g\bigl(f(C_n)\bigr)
= \bigcup_{n \in \N} C_{n+1}
= \bigcup_{n \geq 1} C_n \subseteq C .
\]

\textbf{7.} ليكن $x \neq x'$ في $E$. إذا وقع كلاهما في $C$ فإن $h(x) =
f(x) \neq f(x') = h(x')$ بتباين $f$. وإذا لم يقع أيٌّ منهما في
$C$ فإن $g(h(x)) = x \neq x' = g(h(x'))$، ومنه $h(x) \neq h(x')$. وإذا كان
$x \in C$ و $x' \notin C$ (الحالة المختلطة، بتبديل التسميتين عند
الحاجة): نفترض $h(x) = h(x')$، أي $f(x) = g^{-1}(x')$. وبتطبيق $g$:
$x' = g(f(x)) \in g(f(C))$، ويعطي السؤال 6 أن $x' \in C$ ---
وهذا تناقض. إذن $h(x) \neq h(x')$ في جميع الحالات: أي أن $h$ متباين.

\textbf{8.} ليكن $y \in F$. \emph{الحالة 1: $g(y) \notin C$.} عندئذ
$h(g(y)) = g^{-1}(g(y)) = y$: فالعنصر $g(y)$ سابقة.
\emph{الحالة 2: $g(y) \in C$}، وليكن $g(y) \in C_n$. بما أن $g(y) \in
g(F)$ فإن $g(y) \notin C_0 = E \setminus g(F)$، ومنه $n \geq 1$
و$g(y) \in C_n = g(f(C_{n-1}))$: أي يوجد $x \in C_{n-1}$ يحقق
$g(y) = g(f(x))$. ويعطي تباين $g$ أن $y = f(x)$، ولدينا $x \in
C_{n-1} \subseteq C$، ومنه $h(x) = f(x) = y$. ففي الحالتين يُبلغ $y$:
أي أن $h$ شامل، فهو إذن تقابليّ.

\textbf{9.} إذا كان $E \preceq F$ و $F \preceq E$ فاختر تطبيقين
متباينين $f
\colon E \to F$ و $g \colon F \to E$؛ وتبني الأسئلة 5--8
تقابلًا $h \colon E \to F$، ومنه $E \approx F$. \omterm{def:b1:logic:statement}{والعبارة} غير بديهية
لأن التطبيقين المتباينين المعطيين لا رابط بينهما --- فلا يلزم أن
يكون أيٌّ منهما \omterm{def:b1:logic:inj}{شاملًا}، ولا تعرّف صيغةٌ ساذجة تخلط $f$ و $g$ أيَّ
\omterm{def:b1:logic:map}{تطبيق}: فمضمون المسألة كله هو تجزئة $E$ إلى المنطقة $C$ (حيث ننسخ
$f$) ومتمّمتها (حيث نعكس $g$).

\textbf{10.} (أ) الحقن $\intoo01 \to \intcc01$ متباين؛
\omterm{def:b1:logic:map}{والتطبيق} $x \mapsto \frac{x + 1}3$ يرسل $\intcc01$ تباينًا في
$\intcc{\frac13}{\frac23} \subseteq \intoo01$ (فهو تآلفي ذو معامل توجيه
غير معدوم). وبالسؤال 9 نجد $\intcc01 \approx \intoo01$ --- وهو
تقابل تصعب كتابته صراحةً إلى حد بعيد.
(ب) \emph{التباين.} نفترض $2^p(2q + 1) = 2^{p'}(2q' + 1)$ مع
$p \leq p'$ مثلًا. وبالقسمة على $2^p$: $2q + 1 = 2^{p' - p}(2q' + 1)$.
فإذا كان $p' > p$ كان الطرف الأيمن زوجيًا والطرف الأيسر فرديًا ---
وهذا مستحيل؛ ومنه $p = p'$، ثم $2q + 1 = 2q' + 1$ و $q = q'$.
\emph{الشمول.} نبيّن بالاستقراء القوي أن كل عدد صحيح
$m \geq 1$ على الصورة $2^p(2q + 1)$. من أجل $m = 1$: $p = q = 0$.
وليكن $m \geq 1$ ولنفترض الدعوى صحيحة لكل الأعداد الصحيحة من
$\intint1m$. إذا كان $m + 1$ فرديًا فإن $m + 1 = 2q + 1$ مع $p = 0$.
وإذا كان $m + 1$ زوجيًا فإن $m + 1 = 2m'$ مع $1 \leq m' \leq m$؛
وبالفرض $m' = 2^p(2q + 1)$، ومنه $m + 1 = 2^{p+1}(2q + 1)$. إذن
يبلغ $\varphi(p, q) = 2^p(2q + 1) - 1$ كل $n \in \N$، ويكون
$\varphi$ تقابلًا $\N \times \N \to \N$.

\textbf{11.} بما أن $A$ لا نهائية فإن $A \setminus \{\varphi(0), \dots,
\varphi(n - 1)\}$ لا تكون خالية أبدًا، وخاصية العنصر الأصغر في
$\N$ (المستعملة في البرهان على \cref{thm:b1:logic:induction}) تجعل
التعريف التراجعي مشروعًا. \emph{متزايد تمامًا:}
ينتمي $\varphi(n + 1)$ إلى $A \setminus \{\varphi(0), \dots,
\varphi(n)\} \subseteq A \setminus \{\varphi(0), \dots, \varphi(n -
1)\}$، وأصغر عناصرها هو $\varphi(n)$؛ ومنه $\varphi(n + 1) \geq
\varphi(n)$، والتساوي مستبعد، ومنه $\varphi(n+1) >
\varphi(n)$. \emph{$\varphi(n) \geq n$:} بالاستقراء، $\varphi(0)
\geq 0$ و$\varphi(n + 1) \geq \varphi(n) + 1 \geq n + 1$.
و\emph{التباين} ينتج من الرتابة التامة.
و\emph{الشمول على $A$:} نفترض أن عنصرًا $a \in A$ لا يُبلغ أبدًا.
بما أن $\varphi(a + 1) \geq a + 1 > a$ فإن \omterm{def:b1:logic:sets}{مجموعة} الأعداد $n$ التي تحقق
$\varphi(n) > a$ غير خالية؛ وليكن $n$ أصغر عناصرها. عندئذ يكون
$\varphi(k) \leq a$ لكل $k < n$، ومنه $\varphi(k) < a$ (إذ إن $a$ لا
يُبلغ). فينتمي $a$ عندئذ إلى $A \setminus \{\varphi(0), \dots,
\varphi(n - 1)\}$ ويحقق $a < \varphi(n)$، وهذا يناقض الأصغرية التي
تعرّف $\varphi(n)$. إذن $\varphi$ تقابل $\N \to A$، وتكون
$A$ قابلة للعدّ.

\textbf{12.} ليكن $E \preceq \N$ عبر \omterm{def:b1:logic:inj}{تطبيق متباين} $f$؛ عندئذ $E \approx
f(E)$ (السؤال 2). فإذا كانت $f(E)$ منتهية كانت $E$ منتهية؛ وإذا
كانت $f(E)$ لا نهائية أعطى السؤال 11 أن $f(E) \approx \N$، ومنه
$E \approx \N$ بالتعدّي (السؤال 1). وبالعكس، من الواضح أن المجموعات
المنتهية والمجموعات القابلة للعدّ تُحقن في $\N$. وأمّا الاختصار: فإن
$E \preceq \N$ و$\N
\preceq E$ يعطيان $E \approx \N$ مباشرة بمبرهنة
كانتور--شرودر--برنشتاين --- دون الحاجة إلى أيّ حجة تعداد.

\textbf{13.} ليكن $f \colon E \to \N$ و $g \colon F \to \N$
تطبيقين متباينين. عندئذ يكون $(x, y) \mapsto \varphi\bigl(f(x), g(y)\bigr)$ تطبيقًا
\omterm{def:b1:logic:inj}{متباينًا} $E \times F \to \N$: فإذا تطابقت الصور أعطى تباين
$\varphi$ (السؤال 10) أن $f(x) = f(x')$ و $g(y) = g(y')$،
ومنه $x = x'$ و $y = y'$. وأمّا $\Z \times \N^*$: فكلا العاملين
قابل للعدّ (السؤال 4)، ومنه $\Z \times \N^* \preceq \N$؛ وهي
لا نهائية (إذ تحتوي $\{0\} \times \N^*$)، فهي إذن قابلة للعدّ
بالسؤال 12.

\textbf{14.} لكل عدد ناطق $r$ تمثيل وحيد $r =
p/q$ مع $p \in \Z$ و $q \in \N^*$ والكسر في أبسط صورة
(والوحدانية مبرهن عليها في \cref{ch:b1:arith}؛ ومن أجل $r = 0$ خذ
$0/1$). \omterm{def:b1:logic:map}{والتطبيق} $r \mapsto (p, q)$ متباين عندئذ: فالثنائية
تحدّد $r = p/q$. ومنه $\Q \preceq \Z \times \N^* \preceq \N$
بالسؤال 13. وبما أن $\N \subseteq \Q$ يعطي $\N \preceq \Q$، فإن
السؤال 12 (أو مبرهنة كانتور--شرودر--برنشتاين مباشرة) يبيّن أن $\Q
\approx \N$: أي أن الأعداد الناطقة قابلة للعدّ.

\textbf{15.} من أجل كل $n$ ثبّت تطبيقًا \omterm{def:b1:logic:inj}{متباينًا} $f_n \colon E_n \to \N$.
ومن أجل $x \in \bigcup_n E_n$، ليكن $n(x)$ هو \emph{أصغر} دليل $n$
يحقق $x \in E_n$، ونضع $u(x) = \varphi\bigl(n(x), f_{n(x)}(x)\bigr)
\in \N$. إذا كان $u(x) = u(x')$ أعطى تباين $\varphi$ أن $n(x) =
n(x') = n$ و $f_n(x) = f_n(x')$، ومنه $x = x'$ بتباين
$f_n$. إذن يُحقن الاتحاد في $\N$: فهو قابل للعدّ على الأكثر.

\textbf{16.} لتكن $\Psi(F) = \sum_{i \in F} 2^i$ من أجل $F \subseteq \N$
منتهية (مع $\Psi(\emptyset) = 0$). نفترض $F \neq F'$ وليكن $m$
أكبر عنصر تختلفان عنده، وليكن $m \in F \setminus
F'$ (بتبديل التسميتين عند الحاجة). فالعناصر التي $> m$ تنتمي إلى
المجموعتين معًا أو لا تنتمي إلى أيّ منهما، فتسهم في المجموعين
بالقدر نفسه؛ وبمقارنة إسهامات العناصر التي $\leq m$:
\[
\sum_{i \in F,\, i \leq m} 2^i \geq 2^m
> 2^m - 1 = \sum_{k=0}^{m-1} 2^k
\geq \sum_{i \in F',\, i \leq m} 2^i ,
\]
باستعمال المجموع الهندسي من \cref{exo:b1:logic:4}. ومنه $\Psi(F)
\neq \Psi(F')$: أي أن $\Psi$ متباين وأن \omterm{def:b1:logic:sets}{مجموعة} أجزاء $\N$
المنتهية قابلة للعدّ على الأكثر؛ وهي لا نهائية (إذ تحتوي جميع
المجموعات الأحادية)، فهي إذن قابلة للعدّ.

\textbf{17.} أرسل $A \subseteq \N$ إلى دالتها المميِّزة $\mathbf 1_A
\colon \N \to \{0,1\}$ حيث $\mathbf 1_A(n) = 1$ إذا كان $n \in A$ و $0$
فيما عدا ذلك؛ وأرسل $u \in \{0,1\}^{\N}$ إلى $A_u = \{n \in \N : u(n) =
1\}$. والتطبيقان مقلوبان أحدهما للآخر: إذ $A_{\mathbf 1_A} = A$ و
$\mathbf 1_{A_u} = u$ (تحقق من القيمة عند كل $n$). وحسب
\cref{thm:b1:logic:inverse} يكون كلٌّ منهما تقابلًا: $\mathcal{P}(\N)
\approx \{0,1\}^{\N}$.

\textbf{18.} من أجل كل $n$ لدينا $d(n) = 1 - \Phi(n)(n) \neq \Phi(n)(n)$،
فتختلف المتتاليتان $d$ و $\Phi(n)$ عند الدليل $n$: $d \neq
\Phi(n)$. ومنه لا يكون أيّ $\Phi$ \omterm{def:b1:logic:inj}{شاملًا}، وبالسؤال 3 لا يوجد كذلك أيّ
\omterm{def:b1:logic:inj}{تطبيق متباين} $\{0,1\}^{\N} \to \N$: \omterm{def:b1:logic:sets}{فالمجموعة} $\{0,1\}^{\N}$ غير
قابلة للعدّ على الأكثر. وعبر قاموس السؤال 17 يكون \omterm{def:b1:logic:map}{التطبيق} $\Phi
\colon \N \to \{0,1\}^{\N}$ تطبيقًا $f \colon \N \to
\mathcal{P}(\N)$، وتقابل $d$ المجموعةَ $D = \{n : n
\notin f(n)\}$ (وبالفعل $d(n) = 1 \iff \Phi(n)(n) = 0 \iff n \notin
f(n)$): فالحجة القطرية \emph{هي} برهان كانتور على
\cref{exo:b1:logic:11} من أجل $E = \N$.

\textbf{19.} اكتب $x_n = 0.d_1(n)\,d_2(n)\,d_3(n)\dots$ في صورته
السوية وعرّف $\delta_n = 5$ إذا كان $d_n(n) \neq 5$، و $\delta_n = 6$
إذا كان $d_n(n) = 5$، ثم $x = 0.\delta_1\delta_2\delta_3\dots$ ولا
يستعمل هذا النشر إلا الرقمين $5$ و $6$، فهو لا ينتهي بسلسلة من
الأرقام $9$: أي أنه النشر السويّ لعدد حقيقي $x \in \intco01$.
ومن أجل كل $n$ يختلف رقما $x$ و $x_n$ من الرتبة $n$
(إذ $\delta_n \neq d_n(n)$ بحكم الإنشاء)؛ وبما أن النشور السوية
وحيدة فإن $x \neq x_n$. ومنه لا تستنفد أيّ متتالية $\intco01$:
وبالسؤال 3 مرة أخرى تكون $\intco01$ غير قابلة للعدّ على الأكثر.

\textbf{20.} لدينا $\intco01 \subseteq \R$، فأيّ \omterm{def:b1:logic:inj}{تطبيق متباين}
$\R \to \N$ يتقيّد إلى \omterm{def:b1:logic:inj}{تطبيق متباين} على $\intco01$، وهذا يناقض
السؤال 19: أي أن $\R$ غير قابلة للعدّ. ولو كانت
$\R \setminus \Q$ قابلة للعدّ على الأكثر لكانت
$\R = \Q \cup (\R \setminus \Q)$ اتحادًا لمجموعتين قابلتين
للعدّ على الأكثر، فتكون قابلة للعدّ على الأكثر بالسؤال 15 (بأخذ
$E_0 = \Q$ و$E_n = \R \setminus \Q$ من أجل $n \geq 1$) ---
وهذا تناقض. إذن الأعداد الصمّاء غير قابلة للعدّ. وبدقة أكبر:
داخل $\R$ تكوّن الأعداد الناطقة \omterm{pb:b1:logic:1}{مجموعة قابلة للعدّ} بينما متمّمتها
غير قابلة للعدّ؛ فلا يمكن لأيّ تقابل أن يطابق $\R
\setminus \Q$ مع $\Q$ --- فالأعداد الصمّاء «أكثر» من الناطقة تمامًا،
وإن كانت المجموعتان لا نهائيتين وكثيفتين معًا.

\textbf{21.} العدد $p/q$ (مع $q \neq 0$) جذر لكثير الحدود $qX - p$،
وهو غير معدوم وذو معاملات صحيحة. والعدد $\sqrt 2$ جذر
لكثير الحدود $X^2 - 2$. ومن أجل $x = \sqrt 2 + \sqrt 3$: لدينا $x^2 = 5 + 2\sqrt 6$، ومنه
$x^2 - 5 = 2\sqrt 6$ و $(x^2 - 5)^2 = 24$، أي
\[
x^4 - 10x^2 + 1 = 0 :
\]
فالعدد $\sqrt 2 + \sqrt 3$ جذر لكثير الحدود $X^4 - 10X^2 + 1$.

\textbf{22.} أرسل $P = a_0 + a_1X + \dots + a_nX^n$ (من الدرجة $\leq n$
وبمعاملات صحيحة) إلى $(a_0, \dots, a_n) \in \Z^{n+1}$: وهذا \omterm{def:b1:logic:inj}{التطبيق
متباين}، لأن كثير الحدود محدَّد بمعاملاته. وبالاستقراء على $n$:
\omterm{def:b1:logic:sets}{المجموعة} $\Z^1 = \Z$ قابلة للعدّ (السؤال 4)، \omterm{def:b1:logic:sets}{والمجموعة}
$\Z^{n+2} \approx \Z^{n+1} \times \Z$ قابلة للعدّ على الأكثر بالسؤال 13.
إذن كل \omterm{def:b1:logic:sets}{مجموعة} من كثيرات الحدود الصحيحة المحدودة الدرجة قابلة للعدّ
على الأكثر؛ وهي لا نهائية (إذ تحتوي الثوابت)، فهي إذن قابلة للعدّ
بالسؤال 12.

\textbf{23.} \omterm{def:b1:logic:sets}{مجموعة} كل كثيرات الحدود الصحيحة هي $\bigcup_{n \in
\N} \{P : \deg P \leq n,\ P \text{ ذو معاملات صحيحة}\}$، وهي
اتحاد قابل للعدّ من مجموعات قابلة للعدّ: فهي قابلة للعدّ على الأكثر
بالسؤال 15، ولا نهائية، فهي إذن قابلة للعدّ.

\textbf{24.} من أجل كل كثير حدود صحيح غير معدوم $P$، تكون \omterm{def:b1:logic:sets}{مجموعة}
الجذور $R_P = \{x \in \R : P(x) = 0\}$ منتهية (وفيها $\deg P$ عنصرًا
على الأكثر، وهذا مقبول). وبالسؤال 23 يمكن تعداد كثيرات الحدود
الصحيحة غير المعدومة $P_0, P_1, P_2, \dots$؛ عندئذ تكون $\mathcal{A} =
\bigcup_{n \in \N} R_{P_n}$ اتحادًا قابلًا للعدّ من مجموعات منتهية
(فهي إذن قابلة للعدّ على الأكثر): فهي قابلة للعدّ على الأكثر
بالسؤال 15. وهي تحتوي $\Q$ (السؤال 21)، فهي لا نهائية: أي أن
$\mathcal{A}$ قابلة للعدّ.

\textbf{25.} لو كانت $\R \setminus \mathcal{A}$ قابلة للعدّ على الأكثر
لكانت $\R = \mathcal{A} \cup (\R \setminus \mathcal{A})$ قابلة للعدّ على
الأكثر (السؤال 15)، وهذا يناقض السؤال 20. ومنه توجد أعداد متسامية
بل وتكوّن \omterm{def:b1:logic:sets}{مجموعة} غير قابلة للعدّ، في حين أن الأعداد الجبرية ---
وهي تشمل كل عدد يُبنى من الأعداد الصحيحة بالجذور --- لا تكوّن إلا
هيكلًا قابلًا للعدّ داخل $\R$. وخلاصة معمار المسألة: تضع
الأسئلة 1--3 لغة المقارنة؛ وتتيح مبرهنة كانتور--شرودر--برنشتاين
(الأسئلة 5--9) البرهان على تساوي القوة بتطبيقين متباينين سهلين
بدل تقابل واحد بارع، وقد استُعملت من أجل $\intcc01 \approx \intoo01$،
ومن أجل $\Q$، وفي الجزء 5 كله؛ وأمدّ تقابل الازدواج (السؤال
10) الجداءات والاتحادات القابلة للعدّ (السؤالان 13 و 15)، اللذين
أمدّا بدورهما $\Q$ وكثيرات الحدود الصحيحة و $\mathcal{A}$؛
وأعطت الحجة القطرية (السؤالان 18--19) المتراجحة التامة الوحيدة
$\N \prec \R$ التي تجعل القصة كلها غير بديهية. واستنتاج كانتور
لافت فلسفيًا: فالبرهان لا يبرز أيّ عدد متسامٍ البتة، ومع ذلك
يبيّن أن \emph{كل عدد حقيقي تقريبًا} متسامٍ بمعنى تساوي القوة.
أمّا تسمية عدد متسامٍ بعينه --- مثل $\pi$ أو $\eu$ --- فقد
اقتضت رياضيات مختلفة كل الاختلاف وعقودًا أخرى من العمل.
\end{solution}
